\documentclass[10pt,twocolumn,letterpaper]{article}
\usepackage[utf8]{inputenc}
\usepackage[T1]{fontenc}
\usepackage{amsmath,amssymb}
\usepackage{graphicx}
\usepackage{hyperref}
\usepackage{booktabs}
\usepackage{cite}
\usepackage{xcolor}
\usepackage{enumitem}
\usepackage[margin=0.72in]{geometry}

\title{\textbf{Toward Testable In-Context Reinforcement Learning:\\
Candidate Algorithmic State, Experience Compilation, and\\
Staged Adaptation under Shift}}

\author{Research Proposal}
\date{}

\begin{document}
\maketitle

\begin{abstract}
Recent local literature shows that in-context reinforcement learning (ICRL) can emerge in sequence models, meta-RL agents, and large language models, but the evidence remains strongest for existence rather than mechanism. The local notes in this repository repeatedly expose three limitations: the adaptive computation inside successful models is still mostly implicit; history selection methods remain heuristic or confounded with memory; and robustness under negative evidence and distribution shift is weak. This proposal narrows the agenda to one testable chain of claims. Pillar I asks whether a small set of \emph{candidate algorithmic state variables} can explain and improve in-context adaptation better than opaque hidden state alone. Pillar II turns that hypothesis into a concrete \emph{experience compilation} architecture with explicit objectives, baselines, and ablations for retention and invalidation. Pillar III evaluates the resulting system in a staged manner, moving from controlled offline and procedural-shift settings to bounded embodied validation only if earlier milestones succeed. The goal is not to claim general adaptive intelligence, but to build a grounded framework for testing, implementing, and evaluating interpretable ICRL.
\end{abstract}

\section{Introduction}
The central promise of in-context reinforcement learning is compelling: an agent with fixed weights should improve online by conditioning on interaction history rather than by performing expensive test-time gradient updates. Local evidence in this repository already shows that this promise is not hypothetical. Algorithm Distillation (AD) demonstrates that a causal Transformer can ingest learning histories and improve within new tasks~\cite{laskin2022ad}. Decision-Pretrained Transformers (DPT) show that supervised pretraining can induce online exploration and offline conservatism in bandits and MDPs~\cite{lee2023dpt}. AdA scales black-box meta-RL to open-ended 3D task spaces and exhibits human-timescale adaptation~\cite{adaptive2023ada}. AMAGO-2 further shows that multi-task adaptation is bottlenecked not just by capacity but by return-scale-sensitive optimization, and that classification-style actor--critic objectives can remove a major barrier~\cite{grigsby2024amago2}. Yet the same local notes also reveal why the field is not ready to claim robust adaptive intelligence.

Three recurring weaknesses motivate this proposal. First, successful ICRL systems still behave like \emph{black boxes}. Theory and mechanistic work suggest that Transformers can implement temporal-difference or posterior-sampling-like computations in context~\cite{lin2023decisionmakers,wang2024td,demircan2024sae}, but the internal state of these computations is neither explicit nor directly testable. Second, current systems are still poor at deciding \emph{what from history matters}. Retrieval, filtering, cross-episodic curriculum, and long-context training all help, but the local notes on RA-DT, LHF, and LMAct jointly suggest that current methods still lack a validated mechanism for deciding which parts of raw history should be retained, revised, or discarded for future decisions~\cite{schmied2024radt,chen2025lhf,ruoss2024lmact}. Third, uncertainty handling remains brittle. The local notes on LLM-based in-context exploration repeatedly show difficulty with negative evidence, prompt instability, and exploration collapse, while offline ICRL notes show that better value learning helps more than pure action imitation~\cite{krishnamurthy2024explore,nie2024evolve,monea2024llms,tarasov2025qlearning}.

This proposal therefore turns the strongest three-pillar agenda from the local synthesis into one coherent research program:
\begin{itemize}[leftmargin=1.2em]
\item \textbf{Pillar I: Candidate Algorithmic State for ICRL.} Identify and test explicit latent variables that may support in-context adaptation.
\item \textbf{Pillar II: Experience Compilation and Adaptive Memory.} Build a model that compresses raw history into revisable memory and compare it against strong context baselines.
\item \textbf{Pillar III: Staged Belief-Calibrated Adaptation under Shift.} Validate whether explicit state and compiled memory improve calibration, abstention, and recovery under controlled distribution shift.
\end{itemize}

The agenda is sequential rather than additive. Pillar I defines a candidate object to be tested rather than assumed, Pillar II implements the architecture that can maintain it, and Pillar III validates whether that object is enough to support bounded adaptation under shift.

\section{Literature Review and Research Gaps}
\subsection{Existing evidence for in-context adaptation}
The local note corpus supports the conclusion that in-context adaptation is real across several training paradigms. AD treats an RL algorithm as a long-history-conditioned policy and shows adaptation across adversarial bandits, DarkRoom tasks, and Watermaze, while also explicitly flagging long episodes and multi-episode context as a limitation~\cite{laskin2022ad}. DPT shows that supervised pretraining with optimal actions can produce online exploration and offline decision making in context~\cite{lee2023dpt}. Theoretical work strengthens this picture: \emph{Transformers as Decision Makers} proves that supervised-pretrained Transformers can approximate LinUCB, Thompson sampling, and UCB-VI under suitable assumptions~\cite{lin2023decisionmakers}. Mechanistic work goes further: \emph{Transformers Learn Temporal Difference Methods for In-Context Reinforcement Learning} argues that TD-like evaluation can emerge in the forward pass, and sparse autoencoder analysis finds TD-error-like features in LLM residual streams~\cite{wang2024td,demircan2024sae}.

However, the local notes also show that the field has mostly validated \emph{existence}, not \emph{understanding}. None of the above methods expose a testable internal adaptive state. The corresponding local notes repeatedly emphasize hidden computation, realizability assumptions, or simple environments, which leaves a major gap between ``a model can adapt'' and ``we can identify which internal quantities support that adaptation.''

\subsection{Data and objective design are now central}
The local notes show a clear shift away from assuming clean optimal supervision. DIT replaces optimal labels with weighted pseudo-labels derived from discounted future returns~\cite{dong2024dit}. SAD shows that random-policy data can be useful if learning is restricted to a trust horizon~\cite{dong2024sad}. LHF demonstrates that filtering histories by stability and improvement improves several backbones without changing model structure~\cite{chen2025lhf}. Reward-prediction-based Decision Transformers go further and learn multi-task structured bandits without direct optimal-action labels, but the local note explicitly emphasizes dependence on data diversity: if the demonstrator explores too little, the model cannot infer reward structure well~\cite{mukherjee2024rewardpredictiondt}. Finally, the local note on \emph{Yes, Q-learning Helps Offline In-Context RL} shows that explicit RL objectives improve performance by a large margin over AD-style pure imitation in many offline settings~\cite{tarasov2025qlearning}.

These local notes jointly identify a deeper issue: the field still lacks a principled match between \emph{what the objective optimizes} and \emph{what adaptive behavior requires}. Sequence imitation is often only an indirect proxy for latent quantities that appear to matter in local notes, such as uncertainty, task belief, and delayed credit assignment.

\subsection{Long context alone is not enough}
Several notes in the local corpus make the same point from different angles. RA-DT uses retrieval to reduce the need for full-episode context but still concludes that memory and meta-learning remain confounded in complex environments~\cite{schmied2024radt}. ReLIC scales visual embodied ICRL to 64k steps with partial updates and Sink-KV, but its local note explicitly says that in-context learning only emerges when training diversity prevents simple memorization~\cite{elawady2024relic}. Structured SSMs, Decision Mamba, xLSTM, and FlashAttention-2 all attack sequence efficiency from different levels of the stack~\cite{lu2023ssmrl,huang2024decisionmamba,beck2024xlstm,dao2023flashattention2}. Yet the common pattern in the local notes is that speedups and longer windows solve only the \emph{symptom} of long context. They do not themselves specify how an agent should summarize, revise, or invalidate past experience.

The strongest negative evidence comes from multimodal and LLM settings. LMAct reports that even with many demonstrations and very long contexts, frontier models often fail to turn declarative knowledge into good action policies~\cite{ruoss2024lmact}. The local notes for \emph{Can Large Language Models Explore In-Context?} and EVOLvE both show that raw history is a weak substrate for exploration; models perform much better when history is summarized or scaffolded in a more algorithmic way~\cite{krishnamurthy2024explore,nie2024evolve}. Taken together, these notes motivate testing whether explicit history compilation can outperform raw-history conditioning, not assuming in advance that one specific compiler abstraction is already warranted by the literature.

\subsection{Research gaps motivating the proposal}
The strongest local gaps, already synthesized in the repository, are:
\begin{itemize}[leftmargin=1.2em]
\item lack of a testable notion of algorithmic state for ICRL;
\item inability to separate memory exploitation from genuine online learning;
\item poor handling of uncertainty and negative evidence in both sequence-model and LLM-based ICRL;
\item lack of a validated mechanism for compressing, revising, and invalidating experience summaries;
\item persistent objective mismatch between action imitation and adaptive decision making;
\item weak robustness to variable interfaces, action semantics, and multimodal demonstration formats;
\item evaluation protocols that still underweight abstention, calibration, and non-stationary recovery.
\end{itemize}

The three pillars below are designed to close these gaps in a single, staged agenda.

\section{Three Pillars}
\subsection{Pillar I: Candidate Algorithmic State for ICRL}
\paragraph{Motivation.}
The local note corpus repeatedly implies that ICRL works because a model performs some internal update over history, but the relevant state remains implicit. Local notes on TD-style in-context learning, sparse autoencoder analysis, and partial-MDP inference suggest several \emph{candidate} latent quantities that may matter, including value-related estimates, uncertainty, task belief, and delayed credit~\cite{wang2024td,demircan2024sae,jiang2023partialmdp}. The evidence does not yet justify a universal or minimal decomposition, so the first task is to test which explicit state variables help beyond an opaque hidden state.

\paragraph{Proposed methodology.}
I will study a family of \emph{candidate algorithmic states} rather than a single assumed decomposition. The simplest explicit family augments a standard recurrent or sequence-model latent with heads for value, uncertainty, latent-task belief, and delayed-credit signals. The core research questions are:
\begin{enumerate}[leftmargin=1.3em]
\item Which candidate state variables measurably improve in-context adaptation over a matched opaque latent baseline?
\item Can counterfactual supervision over reward delay, hidden task identity, and uncertainty make these variables causally interpretable?
\item Under which controlled distribution shifts does explicit state improve regret, calibration, or intervention fidelity?
\end{enumerate}
The plan is to start from synthetic and low-cost procedural task families where reward timing, hidden task identity, and uncertainty can be intervened on directly. Matched-capacity baselines will include a raw history-conditioned sequence model, a free-latent variant without typed heads, and objective-level alternatives motivated by DIT, SAD, and conservative offline ICRL~\cite{dong2024dit,dong2024sad,tarasov2025qlearning}. Primary acceptance criteria are improved held-out regret under controlled shifts, lower calibration error, and higher intervention fidelity after counterfactual edits to task identity or reward timing.

\paragraph{Challenges \& Plan B.}
The main risk is that explicit heads simply mirror a useful hidden representation without improving behavior. If strict typed decomposition does not help, the fallback is a semi-structured model: the agent keeps a free latent channel, but must still expose value-, uncertainty-, and belief-related heads that can be evaluated under interventions.

\subsection{Pillar II: Experience Compilation and Adaptive Memory}
\paragraph{Motivation.}
The local notes suggest that the dominant failure mode is not insufficient token budget but insufficient state abstraction. LHF, CEC, RA-DT, and LMAct all show that what gets placed into context matters at least as much as how many tokens are available~\cite{chen2025lhf,shi2023cec,schmied2024radt,ruoss2024lmact}. Therefore, instead of treating trajectories as flat sequences, the next step is to test whether explicit history compilation yields better decisions at the same or lower context budget.

\paragraph{Proposed methodology.}
I propose an \emph{Experience Compiler} that transforms raw histories into event-centric memory updates. The architecture has three components:
\begin{enumerate}[leftmargin=1.3em]
\item a compiler that emits task-relevant event tokens such as hypothesis revisions, failures, surprises, and reward-significant transitions;
\item a revisable memory that can write, retrieve, and invalidate event summaries;
\item a sequence controller that operates on compiled memory rather than raw trajectories.
\end{enumerate}
This system will be trained with a concrete three-part objective: (i) downstream action and value prediction from compiled state, (ii) retention tests that require memory to preserve future decision-relevant information under compression, and (iii) invalidation losses that penalize stale memory when task rules or reward mappings change. Baselines will include a raw-history Transformer, LHF-style filtered context, RA-DT retrieval, and a simple learned summarizer without memory revision~\cite{chen2025lhf,schmied2024radt}. Success means matching or improving decision quality at lower context cost while also passing retention and invalidation ablations. The architecture will be evaluated first in settings where long raw context is known to be inefficient, then only later in bounded multimodal long-demonstration settings where knowing-doing gaps are severe.

\paragraph{Challenges \& Plan B.}
The hardest problem is deciding whether compiled memory has truly retained task-relevant information rather than merely compressing easy trajectories. If fully learned compilation is unstable, a fallback strategy is to bootstrap with weak event supervision from rule changes, sparse rewards, or task switches in procedurally generated worlds, then relax toward self-supervised event discovery.

\subsection{Pillar III: Staged Belief-Calibrated Adaptation under Shift}
\paragraph{Motivation.}
The local notes show that uncertainty handling is a major bottleneck. LLM-based ICRL notes report difficulty using negative episodes and maintaining stable exploration~\cite{monea2024llms,krishnamurthy2024explore,nie2024evolve}. Offline ICRL notes show that better value learning helps more than imitation, but still struggle under poor coverage and trust-horizon limits~\cite{dong2024sad,tarasov2025qlearning}. Embodied notes show that even strong agents rely on diversity, distillation, or curriculum to avoid collapse~\cite{adaptive2023ada,elawady2024relic}. The missing ingredient is a decision rule that uses explicit beliefs and uncertainty to decide whether to exploit, explore, revise memory, or abstain.

\paragraph{Proposed methodology.}
This pillar adds a \emph{belief-calibrated conservative decision layer} on top of Pillars I and II, but in a staged validation sequence rather than a single large deployment jump. The model will output:
\begin{itemize}[leftmargin=1.2em]
\item a conservative value estimate,
\item an epistemic uncertainty estimate,
\item a task-belief summary,
\item an abstention or delay action when evidence is insufficient.
\end{itemize}
Stage 1 will test offline and procedural-shift settings with fixed rules for when abstention is allowed. Stage 2 will add non-stationary interaction and negative-evidence tasks. Stage 3 will attempt a bounded embodied extension only if the earlier stages show gains over conservative offline ICRL and calibration-only baselines. Evaluation will go beyond return to include expected calibration error, regret after rule shifts, abstention-coverage curves, memory-revision frequency, and catastrophic failure rate. Core baselines will include DIT-style relabeling, SAD trust-horizon training, conservative offline ICRL, and calibration-only variants without explicit memory revision~\cite{dong2024dit,dong2024sad,tarasov2025qlearning}.

\paragraph{Challenges \& Plan B.}
The main risk is that calibration objectives destabilize policy learning or encourage over-conservatism. If full abstention-aware learning proves too brittle, the fallback is to stop at calibration-only diagnostics and delayed-action penalties before introducing explicit abstention actions in any embodied setting.

\subsection{Cross-Pillar Evaluation Logic}
To keep the agenda falsifiable, each pillar has a bounded comparison set and explicit failure modes.
\begin{itemize}[leftmargin=1.2em]
\item \textbf{Pillar I fails} if explicit candidate state does not improve any of regret under shift, calibration, or intervention fidelity over a matched-capacity opaque latent.
\item \textbf{Pillar II fails} if compiled memory cannot beat raw-history, filtered, retrieval-based, or simple summarization baselines under either equal context budget or equal compute budget.
\item \textbf{Pillar III fails} if calibration or abstention reduces catastrophic failures only by collapsing coverage or reward beyond a preset operating frontier.
\end{itemize}
This staged logic is intended to prevent the proposal from hiding weak mechanism under larger models, longer context, or heavier infrastructure.

\section{Expected Outcomes}
If successful, this proposal will contribute:
\begin{itemize}[leftmargin=1.2em]
\item an evidence-grounded test bed for candidate algorithmic state in ICRL rather than relying only on black-box hidden-state stories;
\item a concrete context-compilation architecture with measurable retention and invalidation tests;
\item a staged protocol for uncertainty-aware adaptation under coverage limits and distribution shift;
\item a bounded evaluation recipe that makes failure under non-stationarity, calibration, and abstention visible.
\end{itemize}

The expected research outputs include:
\begin{enumerate}[leftmargin=1.3em]
\item a mechanism paper on candidate algorithmic state and intervention-based tests,
\item a methods paper on experience compilation and adaptive memory,
\item a validation paper on calibration and adaptation under controlled shift, with embodied extension only if the earlier stages succeed.
\end{enumerate}

\section{Timeline}
\begin{table}[h]
\centering
\small
\begin{tabular}{p{0.15\linewidth}p{0.76\linewidth}}
\toprule
\textbf{Phase} & \textbf{Plan} \\
\midrule
Months 1--4 & Reproduce note-backed baselines, formalize candidate-state hypotheses, and build intervention-friendly synthetic tasks. \\
Months 5--10 & Develop and test candidate-state models; run counterfactual credit, belief, and uncertainty experiments; submit a mechanism paper. \\
Months 11--18 & Build the experience compiler and revisable memory architecture; benchmark against AD, RA-DT, LHF, summarization, and long-context baselines. \\
Months 19--26 & Add belief-calibrated conservative adaptation; evaluate offline and non-stationary settings with fixed abstention rules; submit a methods paper. \\
Months 27--32 & Attempt a bounded embodied or multimodal extension only if the earlier stages meet the acceptance criteria; otherwise deepen controlled-shift evaluation. \\
Months 33--36 & Consolidate results, write dissertation chapters, and release code, artifacts, and negative-result documentation where needed. \\
\bottomrule
\end{tabular}
\caption{Proposed research timeline.}
\end{table}

\bibliographystyle{plain}
\bibliography{references}

\end{document}
